\chapter{§1.1 整数的认识}
\label{sec:1.1}


\gradelevel{1-3 年级}


\section*{自然数与 0}


\begin{explore}


开学了，老师让大家数一数教室里有多少把椅子。你一把一把地指着数：1、2、3……一直数到 40。但放学后，教室里一个人也没有了——这时候教室里有多少个学生呢？

\end{explore}

\begin{understand}


当我们数物体的个数时，用到的数 $1, 2, 3, 4, 5, \ldots$ 叫做\textbf{自然数}。自然数从 $1$ 开始，可以一直往后数，没有尽头。

当一个物体也没有的时候，我们用 \textbf{ $0$ } 来表示。 $0$ 也是自然数。

所以，自然数就是：

\[
0, 1, 2, 3, 4, 5, 6, 7, 8, 9, 10, 11, \ldots
\]


每一个自然数都比它前面那个大 $1$ 。最小的自然数是 $0$ ，没有最大的自然数。

\end{understand}

\begin{workedexamples}


\textbf{例 1.} 从下面的数中选出自然数： $3$ ， $0$ ， $-2$ ， $5$ ， $1.5$ 。

\textbf{解：} 自然数是 $0, 1, 2, 3, \ldots$ 这些数。
\begin{itemize}
  \item $3$ 是自然数。
  \item $0$ 是自然数。
  \item $-2$ 是负数，不是自然数。
  \item $5$ 是自然数。
  \item $1.5$ 是小数，不是自然数。
\end{itemize}


所以自然数有： $0$ 、 $3$ 、 $5$ 。

\textbf{例 2.} 在 $0$ 、 $6$ 、 $11$ 、 $100$ 这四个数中，最小的是哪个？最大的是哪个？

\textbf{解：} 比较大小： $0 < 6 < 11 < 100$ 。最小的是 $0$ ，最大的是 $100$ 。

\end{workedexamples}

\begin{keytakeaway}


\begin{itemize}
  \item 自然数包括 $0$ 和所有正整数 $1, 2, 3, \ldots$
  \item 最小的自然数是 $0$ ，没有最大的自然数。
  \item $0$ 表示"一个也没有"。
\end{itemize}


\end{keytakeaway}

\begin{practice}


\begin{enumerate}
  \item 写出 $5$ 以内的所有自然数。
  \item 下面哪些是自然数？ $7$ ， $0$ ， $-1$ ， $2.3$ ， $15$ 。
  \item 按从小到大的顺序排列： $9$ ， $0$ ， $4$ ， $12$ ， $1$ 。
\end{enumerate}


\end{practice}

\section*{计数与数位}


\begin{explore}


妈妈给你 10 颗糖，你又得到了 3 颗。现在你有 13 颗糖。"13"这个数里面的"1"和"3"表示的意思一样吗？

\end{explore}

\begin{understand}


我们用十个数字符号来写所有的数： $0, 1, 2, 3, 4, 5, 6, 7, 8, 9$ 。

同一个数字放在不同的位置上，表示不同的大小。这就是\textbf{数位}的概念。

以 $13$ 为例：
\begin{itemize}
  \item "3"在\textbf{个位}上，表示 $3$ 个一。
  \item "1"在\textbf{十位}上，表示 $1$ 个十。
\end{itemize}


所以 $13 = 1 \times 10 + 3 \times 1$ 。

数位从右到左依次是：\textbf{个位、十位、百位、千位}……每相邻两个数位之间是\textbf{十进关系}——满十进一。

\begin{center}
\begin{tabular}{llll}
\toprule
千位 & 百位 & 十位 & 个位 \\
\midrule
 &  & 1 & 3 \\
\bottomrule
\end{tabular}
\end{center}


再看 $205$ ：
\begin{itemize}
  \item 个位上是 $5$ ，表示 $5$ 个一。
  \item 十位上是 $0$ ，表示 $0$ 个十。
  \item 百位上是 $2$ ，表示 $2$ 个百。
\end{itemize}


所以 $205 = 2 \times 100 + 0 \times 10 + 5 \times 1$ 。

\end{understand}

\begin{workedexamples}


\textbf{例 1.} 写出 $347$ 各数位上的数字及其表示的意义。

\textbf{解：}
\begin{itemize}
  \item 个位上是 $7$ ，表示 $7$ 个一。
  \item 十位上是 $4$ ，表示 $4$ 个十。
  \item 百位上是 $3$ ，表示 $3$ 个百。
\end{itemize}


$347 = 3 \times 100 + 4 \times 10 + 7 \times 1$ 。

\textbf{例 2.} 用 $6$ 、 $0$ 、 $2$ 三个数字，组成一个最大的三位数和一个最小的三位数。

\textbf{解：}
\begin{itemize}
  \item 最大的三位数：把最大的数字放在最高位， $620$ 。
  \item 最小的三位数：把最小的非零数字放在最高位（百位不能是 $0$ ）， $206$ 。
\end{itemize}


\end{workedexamples}

\begin{keytakeaway}


\begin{itemize}
  \item 我们使用"十进制"：满十向前一位进一。
  \item 同一个数字在不同数位上表示不同的大小。
  \item 读数从高位到低位，中间有零要读出。
\end{itemize}


\end{keytakeaway}

\begin{practice}


\begin{enumerate}
  \item $586$ 的百位上是几？表示什么？
  \item 用 $3$ 、 $7$ 、 $0$ 三个数字组成最大的三位数和最小的三位数。
  \item 写出 $4019$ 各数位上的数字及其意义。
\end{enumerate}


\end{practice}

\section*{万以内数的认识与大小比较}


\begin{explore}


学校图书馆有 $3860$ 本书，隔壁学校有 $4012$ 本。哪个学校的书多？你是怎么比较的？

\end{explore}

\begin{understand}


当数扩展到千位以上，我们需要认识\textbf{千位}和\textbf{万位}。

\[
\text{万位} \quad \text{千位} \quad \text{百位} \quad \text{十位} \quad \text{个位}
\]


例如， $3860$ 读作"三千八百六十"。

\textbf{大小比较}的方法：
\begin{enumerate}
  \item \textbf{位数不同}——位数多的数更大。例如 $986 < 1002$ ，三位数一定比四位数小。
  \item \textbf{位数相同}——从最高位开始逐位比较。例如比较 $3860$ 和 $4012$ ：千位上 $3 < 4$ ，所以 $3860 < 4012$ 。
\end{enumerate}


\end{understand}

\begin{workedexamples}


\textbf{例 1.} 比较 $5678$ 和 $5702$ 的大小。

\textbf{解：} 两个数都是四位数，从最高位开始比较。
\begin{itemize}
  \item 千位：都是 $5$ ，相同。
  \item 百位： $6 < 7$ 。
\end{itemize}


所以 $5678 < 5702$ 。

\textbf{例 2.} 把 $320$ 、 $3200$ 、 $302$ 、 $3020$ 按从小到大的顺序排列。

\textbf{解：} 先看位数： $320$ 和 $302$ 是三位数， $3200$ 和 $3020$ 是四位数。三位数一定比四位数小。
\begin{itemize}
  \item 三位数之间： $302 < 320$ （百位相同，十位 $0 < 2$ ）。
  \item 四位数之间： $3020 < 3200$ （千位相同，百位 $0 < 2$ ）。
\end{itemize}


从小到大： $302 < 320 < 3020 < 3200$ 。

\end{workedexamples}

\begin{keytakeaway}


\begin{itemize}
  \item 万以内的数位从右到左依次是：个位、十位、百位、千位、万位。
  \item 比较大小：先看位数，位数多的大；位数相同从最高位起逐位比较。
  \item 数的读写要注意中间和末尾的零。
\end{itemize}


\end{keytakeaway}

\begin{practice}


\begin{enumerate}
  \item 用 $>$ 、 $<$ 或 $=$ 填空： $4500 \enspace\underline{\quad}\enspace 4050$ ， $999 \enspace\underline{\quad}\enspace 1000$ 。
  \item 把 $6010$ 、 $601$ 、 $6100$ 、 $610$ 从小到大排列。
  \item 一个四位数，千位上是 $2$ ，百位上是 $0$ ，十位上是 $5$ ，个位上是 $8$ ，这个数是多少？读作什么？
\end{enumerate}


\end{practice}

\section*{整数的分类}


\begin{explore}


在温度计上，零度以上的温度用正数表示，比如 $+5°\text{C}$ （零上 $5$ 度）。零度就是 $0°\text{C}$ 。那零度以下呢？到了初中你会学习负数（§2.1），但在小学阶段，我们先认识正整数和 $0$ 。

\end{explore}

\begin{understand}


在小学阶段我们打交道的整数可以分为两类：

\begin{itemize}
  \item \textbf{正整数}： $1, 2, 3, 4, 5, \ldots$ ——大于零的整数。
  \item \textbf{零}： $0$ ——既不是正整数，也不是负整数。
\end{itemize}


正整数和 $0$ 合在一起，就是我们熟悉的\textbf{自然数}。

\end{understand}

\begin{quote}
到了七年级（§2.1 有理数），你还会认识\textbf{负整数}，如 $-1, -2, -3, \ldots$ ，到那时整数的家族就完整了。
\end{quote}


\begin{workedexamples}


\textbf{例 1.} 把下面的数分成两组：正整数和零。 $8$ ， $0$ ， $15$ ， $100$ ， $0$ ， $3$ 。

\textbf{解：}
\begin{itemize}
  \item 正整数： $8$ ， $15$ ， $100$ ， $3$ 。
  \item 零： $0$ 。
\end{itemize}


\textbf{例 2.} 最小的正整数是几？有没有最大的正整数？

\textbf{解：} 最小的正整数是 $1$ 。没有最大的正整数，因为任何正整数加 $1$ 都能得到一个更大的正整数。

\end{workedexamples}

\begin{keytakeaway}


\begin{itemize}
  \item 正整数就是 $1, 2, 3, \ldots$ ，大于 $0$ 。
  \item $0$ 是自然数，但不是正整数。
  \item 自然数 $=$ 正整数 $+$ 零。
\end{itemize}


\end{keytakeaway}

\begin{practice}


\begin{enumerate}
  \item 最小的自然数是多少？最小的正整数是多少？
  \item 判断对错： $0$ 是正整数。
  \item 写出大于 $5$ 且小于 $10$ 的所有正整数。
\end{enumerate}


\end{practice}



\begin{answer}


\textbf{练一练 自然数与 0}

\begin{enumerate}
  \item $0$ 、 $1$ 、 $2$ 、 $3$ 、 $4$ 、 $5$ 。
  \item 自然数有： $7$ 、 $0$ 、 $15$ 。（ $-1$ 是负数， $2.3$ 是小数，都不是自然数。）
  \item $0 < 1 < 4 < 9 < 12$ 。
\end{enumerate}


\textbf{练一练 计数与数位}

\begin{enumerate}
  \item 百位上是 $5$ ，表示 $5$ 个百。
  \item 最大的三位数： $730$ ；最小的三位数： $307$ 。
  \item 千位上是 $4$ ，表示 $4$ 个千；百位上是 $0$ ，表示 $0$ 个百；十位上是 $1$ ，表示 $1$ 个十；个位上是 $9$ ，表示 $9$ 个一。
\end{enumerate}


\textbf{练一练 万以内数的认识与大小比较}

\begin{enumerate}
  \item $4500 > 4050$ ， $999 < 1000$ 。
  \item $601 < 610 < 6010 < 6100$ 。
  \item 这个数是 $2058$ ，读作"两千零五十八"。
\end{enumerate}


\textbf{练一练 整数的分类}

\begin{enumerate}
  \item 最小的自然数是 $0$ ，最小的正整数是 $1$ 。
  \item 错。 $0$ 不是正整数。
  \item $6$ 、 $7$ 、 $8$ 、 $9$ 。
\end{enumerate}


\end{answer}
